\subsection*{الف}
$\frac{1}{t}\sin(3\pi t)$ معادل $\text{sinc}$ با باند $3\pi$ است؛ از این رو بالاترین فرکانس‌ها:
\begin{flalign*}
    \frac{1}{t}\sin(3\pi t) &\longleftrightarrow \pi\, \text{rect}\!\left(\frac{\omega}{6\pi}\right), \quad |\omega|\le 3\pi \\
    \cos(2\pi t) &\longleftrightarrow \omega = \pm 2\pi \\
    \omega_{max} &= 3\pi \implies \omega_s > 2\omega_{max} = 6\pi \\
    T_{max} &= \frac{2\pi}{\omega_s} = \frac{2\pi}{6\pi} = \frac{1}{3} &&
\end{flalign*}

\subsection*{ب}
ضرب در حوزهٔ زمان معادل کانولوشن در حوزهٔ فرکانس و جمع پهنای باندهاست:
\begin{flalign*}
    \cos(12\pi t) &\longleftrightarrow \omega = \pm 12\pi \\
    \frac{\sin(\pi t)}{2t} &= \frac{\pi}{2}\,\frac{\sin(\pi t)}{\pi t} \longleftrightarrow \text{band} = \pi \\
    \omega_{max} &= 12\pi + \pi = 13\pi \\
    \omega_s &> 26\pi \implies T_{max} = \frac{2\pi}{26\pi} = \frac{1}{13} &&
\end{flalign*}

\subsection*{ج}
کانولوشن در زمان معادل ضرب در فرکانس است؛ پهنای باند برابر حداقل دو طیف می‌شود:
\begin{flalign*}
    e^{-6t}u(t) &\longleftrightarrow \frac{1}{6+j\omega} \quad (\text{band} \to \infty) \\
    \frac{\sin(\omega_a t)}{\pi t} &\longleftrightarrow \text{rect},\ |\omega|\le \omega_a \\
    X_3(j\omega) &= \frac{1}{6+j\omega}\cdot \text{rect}\!\left(\frac{\omega}{2\omega_a}\right) \implies \omega_{max} = \omega_a \\
    \omega_s &> 2\omega_a \implies T_{max} = \frac{2\pi}{2\omega_a} = \frac{\pi}{\omega_a} &&
\end{flalign*}
($\omega$ در صورت سؤال همان $\omega_a$ در نظر گرفته شده است.)

\subsection*{د}
$x_4(t)=w(t)z(t)$، پس $X_4 = \frac{1}{2\pi} W * Z$ و پهنای باند برابر مجموع است:
\begin{flalign*}
    Z(j\omega) &: \ |\omega| \le 4\pi \\
    W(j\omega) &: \ |\omega| \le W_a \\
    \omega_{max} &= 4\pi + W_a \\
    \omega_s &> 2(4\pi + W_a) \implies T_{max} = \frac{2\pi}{2(4\pi+W_a)} = \frac{\pi}{4\pi + W_a} &&
\end{flalign*}